\section{Respuesta a la RQ4}
\label{sec:comparacion:respuesta}

La RQ4 pregunta cómo pueden compararse las APIs de correo de forma sistemática para responder a preguntas de decisión que involucran simultáneamente requisitos funcionales y no funcionales.
La respuesta es un procedimiento que instancia cada una de las 32 plantillas del catálogo sobre cada API.
Cada instancia se responde con el instrumento que declara su clase: un campo de la Datasheet, una operación de PRIME4API, una composición de HARVEY4API o la matriz de \textit{features} por plan.
Cada respuesta queda registrada con la fuente que la sostiene.

Con ese instrumental, 25 de las 32 preguntas tienen una clase que las declara respondibles con la Datasheet, PRIME4API o HARVEY4API.
Si se descuentan las cuatro de HARVEY4API, que no se ha ejecutado, quedan 21.
Otras 2, fuera de esas 25, se responden leyendo la matriz de \textit{features}.
Las cinco restantes señalan lo que falta, como la concurrencia de peticiones, que ACM2YAML v0.3 no serializa, o la retención de los registros de envío, un límite que no se gasta con cada envío.

Cada API responde entre 18 de esas 25 preguntas, en Postmark, y 22, en Resend (Tabla~\ref{tab:cobertura-pi}).
En SES, 17 de sus 21 respuestas no se dan sobre el plan de suscripción que presupone la pregunta: 16 se dan sobre el entorno de pruebas y una sobre la tarifa sin suscripción.
En Mailchimp, 10 de 20 son de la cuenta de prueba.

La comparación es asimétrica.
El coste, la cuota y el caudal se comparan calculando sobre la Datasheet con PRIME4API, y los \textit{webhooks} o la IP dedicada se comparan leyendo a mano la matriz de \textit{features}.
El caso de los 100.000 correos al mes necesita las dos: el precio, el caudal y el lote salen de la Datasheet, y los \textit{webhooks}, de la matriz.
Cumplen SendGrid, Mailgun, Postmark y Brevo, y la más barata es SendGrid, a 34,95 USD al mes con Essentials 100K.
Con los campos de v0.3 saldría a 84,95 USD, 2,4 veces el precio que SendGrid publica para ese volumen, porque la cuota no guarda el volumen de 100.000 correos (Sección~\ref{sec:comparacion:caso}).

La respuesta tiene dos límites.
La base son ocho APIs, y en cinco de ellas lo observado con cuenta propia corresponde a un estado de la suscripción, la fase en que el proveedor sitúa la cuenta, y no a un plan.
Son el entorno de pruebas en SES, la cuenta de prueba en Mailchimp y una cuenta gratuita o sin aprobar en MailerSend, Mailgun y Brevo.
El segundo límite es que las cuatro preguntas que combinan varias Datasheets o varias operaciones cuentan como respondibles por la especificación de HARVEY4API que publica Ruiz López \citep{ruizlopez2026tfm}, sin que la herramienta se haya ejecutado.
Las respuestas que dio Claude Opus 5.5 leyendo las ocho Datasheets llegan a la misma conclusión que la referencia calculada a mano, pero no sustituyen a la herramienta (Sección~\ref{sec:comparacion:harvey}).
